% 2026-11-08 v7. Integration fragment for Japanese LuaLaTeX.
% The new comparison argument is proved below for the explicitly stated
% marked cubic Z-algebra presentation. See 01-research-notes.md for scope.
% Existing bibliography keys are reused. Add QuadricBR2020 and QuadricMU2021.

\subsection{楕円型 quadric の一点 blowup と切断環の比較}
\label{blowup:quadric-comparison}

以下では $k$ を標数零の代数閉体とする。
$A$ を三次元 cubic AS-regular $\mathbb Z$-代数とし、
$X=\operatorname{Proj}A$、$P_i=\pi(e_iA)$ と書く。
$A$ の楕円データを $(E,M_0,M_1,M_2)$ とし、$E$ は滑らかと仮定する。
この記法では $M_i$ は次数2の直線束であり、
$E\hookrightarrow\mathbb P^1\times\mathbb P^1$ を定める隣接する二つの pencil を持つ。
楕円曲線への包含を $i\colon E\hookrightarrow X$ とする。
標準的な同値 $\operatorname{qgr}(A/(g_i))\simeq\operatorname{coh}E$ を
$Li^*P_0=\mathcal O_E$ となるよう規格化すると
\begin{equation}\label{blowup:quadric-restrictions}
 F_0=\mathcal O_E,\qquad
 F_j=Li^*P_j=(M_0\otimes\cdots\otimes M_{j-1})^{-1}
 \quad(1\le j\le3).
\end{equation}
この同定は twisted homogeneous coordinate $\mathbb Z$-代数の定義から従い、
$\operatorname{Hom}(P_{j+1},P_j)=H^0(E,M_j)$ の制限は切断の評価写像である。
用いている cubic $\mathbb Z$-代数の分類と楕円データについては
\cite[命題5.1.2、\S5.5]{VdBQuadrics} を参照する。

\begin{lemma}[一点上の半単純な制限による消滅]
\label{blowup:quadric-kernel}
$T\in D^-_{\mathrm{coh}}(X)$ と $t\in E$ が
\[
 R\Gamma(X,T)=0,\qquad
 \mathcal H^j(Li^*T)\cong k(t)^{\oplus m_j}
 \quad(j\in\mathbb Z)
\]
を満たすとする。各 $m_j$ は有限とする。このとき $T=0$ である。
\end{lemma}

\begin{proof}
$(P_3,P_2,P_1,P_0)$ は full strong exceptional collection である。
最初の二対象を right mutation して
\[
 0\longrightarrow P_3\longrightarrow
 H^0(E,M_2)^\vee\otimes P_2\longrightarrow U\longrightarrow0
\]
を得る。Belmans--Raedschelders の計算
\cite[命題3.12]{QuadricBR2020} により、
$(P_2,U,P_1,P_0)$ は full strong exceptional collection であり、
$U$ と $P_1$ は完全に直交する。
従って $R\Gamma(X,-)$ の核は
\[
 \mathcal K=\langle P_2,U,P_1\rangle
\]
に対応する。
その tilting 代数 $H$ は、一つの頂点から他の二頂点へそれぞれ二本の矢が出る
関係式のない quiver の道代数であり、hereditary である。
ここで必要となるのは三頂点の部分代数だけである。

上の mutation 列を $E$ に制限する。
$M_2$ の二次元の全切断空間は basepoint-free なので、
制限した左の射は部分直線束の包含になる。
二次元ベクトル空間の行列式の基底を一つ固定すれば
\begin{equation}\label{blowup:quadric-mutated-restriction}
 Li^*P_2=F_2,\qquad Li^*U=F_2\otimes M_2,\qquad
 Li^*P_1=F_2\otimes M_1.
\end{equation}
二組の矢は、それぞれ $M_2$ と $M_1$ の全切断による写像である。

hereditary 性より、$D^-_{\mathrm{fd}}(H)$ の対象はそのコホモロジー加群の
シフトの直和に分解する。有限次元 $H$-加群の最小射影分解を
$\mathcal K$ に戻すと、その対応物は
\begin{equation}\label{blowup:quadric-two-term}
 W\otimes P_2\longrightarrow
 V\otimes U\oplus V'\otimes P_1
\end{equation}
という二項複体と $P_2$ の直和に分かれる。
最小性により、始域の $P_2$ から終域の $P_2$ への成分は零になる。
$Li^*P_2$ は直線束なので、$P_2$ の単独の直和因子は仮定から生じない。
また、各二項複体の制限のコホモロジーも $k(t)$ の直和である。
有限振幅の制限関手はこの各次数で有限な直和と可換する。

式\eqref{blowup:quadric-two-term}を制限し $F_2^{-1}$ を掛けると
\[
 \varphi\colon\mathcal O_E^{\oplus w}\longrightarrow
 M_2^{\oplus v}\oplus M_1^{\oplus v'}
\]
を得る。核は局所自由層の部分層かつ torsion なので零である。
余核が torsion であることから $w=v+v'$ である。
直線束の次数を比較すると
\[
 \operatorname{length}(\operatorname{coker}\varphi)=2(v+v')=2w.
\]
一方、余核が $k(t)^{\oplus m}$ なら、$t$ におけるファイバーへの全射から
$m\le v+v'=w$ となる。従って $2w\le w$ であり $w=0$ である。
全ての直和因子が零となり、$T=0$ を得る。
\end{proof}

\begin{theorem}[cubic 表示を持つ quadric の比較同値]
\label{blowup:quadric-section-equivalence}
上の $X$ が、構造層と $E$ の埋め込みを保つ同値
\[
 X\simeq\operatorname{Proj}Q
\]
によって、$A$ の二周期的な even Veronese の通常の斉次表示 $Q$ を持つとする。
さらに $Q$ は domain で、次数2の中心的正則元 $g$ を持ち
\[
 Q/gQ=B(E,\mathcal L,\alpha),\qquad \deg\mathcal L=4
\]
とする。$T_0=Q^{(2)}$ とし、
\[
 T_0/gT_0=B(E,\mathcal M,\tau),\qquad
 \mathcal M=\mathcal L\otimes\alpha^*\mathcal L,\qquad
 \tau=\alpha^2
\]
と書く。$|\tau|=\infty$ と仮定する。
任意の $t\in E$ に対して $R=T_0(t)$ とおけば、
Van den Bergh の一点 blowup $\alpha_t\colon\widetilde X_t\to X$ と切断環への射
$\delta_t\colon\widetilde X_t\to\operatorname{Proj}R$ は
\[
 R\delta_{t*}\colon D^b(\operatorname{coh}\widetilde X_t)
       \xrightarrow{\ \sim\ }D^b(\operatorname{qgr}R)
\]
という dg-enhanced 同値を与える。その逆は $L\delta_t^*$ である。
従って $\mathcal E_t=R\delta_{t*}\mathcal O_{L_t}(-1)$ とおけば
\begin{equation}\label{blowup:quadric-unconditional-sod}
 \operatorname{Perf}(\operatorname{qgr}R)
 =\left\langle
     \mathcal E_t\otimes_k\operatorname{Perf}(k),\quad
     R\delta_{t*}L\alpha_t^*\operatorname{Perf}(\operatorname{qgr}Q)
   \right\rangle
\end{equation}
\end{theorem}

\begin{proof}
Van den Bergh の構成に用いる仮定を確認する。
$E\subset X$ は滑らかな両側閉 Cartier divisor であり、
$\mathcal O_X(E)=\mathcal O_X(2)$ は ample である。
標準的な $P_i$ は $E$-flat で、制限は直線束であり、$X$ を生成する。
cubic AS-regular 性から $\operatorname{Mod}(X)$ は有限 homological dimension を持つ。
また $P_0=\mathcal O_X$ が exceptional なので
$R\Gamma(X,\mathcal O_X)=k$ である。
normal bimodule に点イデアルを掛けた $E$ 上の invertible bimodule は
$(\mathcal M(-t))_\tau$ であり、次数は7である。
その正の冪は $H^1=0$ を満たし ample である。
従って \cite[\S9.2、命題9.2.2、命題9.2.3]{BlowupVdB2001} が適用される。

切断環の同定もここで確認できる。
Van den Bergh の Rees bimodule algebra を $\mathcal D_t$ とし
$D_t=\Gamma(X,\mathcal D_t\mathcal O_X)$ とする。
この標準表示では
$Q_n=H^0(X,\mathcal O_X(n))$ for $n\ge0$ である。
従って Rees algebra の包含から $D_t\subset T_0$ を得る。
同文献の式(6.16)と命題9.2.3から
\[
 0\longrightarrow D_t(-1)\xrightarrow{g}D_t
 \longrightarrow B(E,\mathcal M(-t),\tau)\longrightarrow0
\]
を得る。次数1部分は
\[
 (D_t)_1/k g=H^0(E,\mathcal M(-t))\subset H^0(E,\mathcal M).
\]
この式は点の表示を固定する。
右加群の grading shift による点の移動はこの標識に含める。
$R=T_0(t)$ はこの次数1部分で生成される
\cite[定理5.3]{BlowupRSS2015}。
従って $R\subset D_t$ である。
両者の $g$ による商は同じ twisted homogeneous coordinate ring であり、
各次数の次元は
\[
 \dim_k R_n=\dim_k(D_t)_n=1+\frac{7n(n+1)}2.
\]
よって $R=D_t$ である。

命題9.2.3は $R\delta_{t*}L\delta_t^*=\mathrm{id}$ を $D^-$ 上で与える。
さらに $\delta_{t*}$ の cohomological dimension は高々1である。
$N\in D^-_{\mathrm{coh}}(\widetilde X_t)$ が $R\delta_{t*}N=0$ を満たすとする。
$T=R\alpha_{t*}N$ とおく。
blowup の分解 \cite[定理8.4.1]{BlowupVdB2001} により三角
\[
 L\alpha_t^*T\longrightarrow N\longrightarrow C\longrightarrow
\]
が存在し、$C$ は $\mathcal O_{L_t}(-1)$ のシフトの直和である。
pushforward の合成と $R\delta_{t*}N=0$ から
$R\Gamma(X,T)=R\Gamma(\widetilde X_t,N)=0$ となる。
同文献の式(9.11)と、$\delta_t$ が $E$ 上で恒等写像になることから
$Li^*N=0$ である。
$\mathcal O_{L_t}(-1)$ の $E$ への導来制限は一点の skyscraper である。
従って $Li^*T$ の各コホモロジーは同じ一点の skyscraper の有限直和であり、
補題\ref{blowup:quadric-kernel}により $T=0$ となる。
三角から $N\simeq C$ であり、その制限が零であることから $C=0$、従って $N=0$ である。
これで $R\delta_{t*}$ の核が $D^-_{\mathrm{coh}}$ 上で消滅した。

有界性は次の spectral sequence から従う。
$F\in\operatorname{coh}(\operatorname{Proj}R)$ とし $N=L\delta_t^*F$ とおく。
$R^j\delta_{t*}=0$ for $j>1$ より
\[
 0\longrightarrow R^1\delta_{t*}\mathcal H^{r-1}(N)
 \longrightarrow\mathcal H^r(R\delta_{t*}N)
 \longrightarrow\delta_{t*}\mathcal H^r(N)\longrightarrow0
\]
を得る。$R\delta_{t*}N=F$ なので、$r<-1$ に対して
$R\delta_{t*}\mathcal H^r(N)=0$ である。
核の消滅から $\mathcal H^r(N)=0$ となり、$L\delta_t^*$ の振幅は $[-1,0]$ に入る。
従って上の随伴同値は $D^b_{\mathrm{coh}}$ に制限する。
導来関手は元の bimodule functor から構成されるので dg enhancement と両立する。
最後に定理8.4.1の分解をこの同値で移すと
式\eqref{blowup:quadric-unconditional-sod}を得る。
\end{proof}

\begin{remark}[適用範囲と generic localization]
この定理の仮定は、$E$ と構造層を標識した cubic $\mathbb Z$-代数表示を含む。
この表示を既に持つ曲面族では、categorical blowup と切断環の比較を
さらに仮定する必要がなくなる。
RSS の任意の hypersurface 表示の quadric に適用する際には、
その標識付き表示との同定を確認する。

点イデアル bimodule は generic localization で単位 bimodule になる。
従って Rees bimodule algebra の局所化は一変数の斉次環となり、
$\alpha_t$ は generic quotient 上で恒等同値を誘導する。
また $D_t=R\subset T_0$ という上の同定は、$\delta_t$ が誘導する
generic quotient の同一視を、この包含が誘導する分数体の同一視に固定する。
従って、$D=Q_0(Q)\cong Q_0(R)$ とし両側の局所化関手を $\ell_Q,\ell_R$ と書けば
\[
 \ell_R\circ R\delta_{t*}L\alpha_t^*\simeq\ell_Q
\]
となる。例外対象 $\mathcal E_t$ の generic rank は零である。
ここで使う $Q_0(R)=Q_0(T_0)$ は
\cite[定理1.1(2)]{BlowupRSS2015}の分数環の保存とも一致する。
\end{remark}

\begin{proposition}[Sklyanin pencil の一般の quadric への適用]
\label{blowup:generic-sklyanin-quadric}
$S_4$ を、滑らかな楕円曲線 $E$、次数4の直線束 $\mathcal L$、
無限位数の平行移動 $\alpha$ に付随する四次元 Sklyanin 代数とする。
その中心的二次元部分空間から得られる pencil
$Q=S_4/(\Omega)$ の非空 Zariski 開集合上で、任意の $t\in E$ に対し
定理\ref{blowup:quadric-section-equivalence}の比較同値と
式\eqref{blowup:quadric-unconditional-sod}が成立する。
\end{proposition}

\begin{proof}
滑らかな $\operatorname{qgr}Q$ における二つの正規化された rank-one
maximal Cohen--Macaulay module を $X,Y$ とする。
各々は次数0の二元で生成される。
syzygy による置換 $M\mapsto\Omega M(1)$ が $X,Y$ を交換するとき、
$Q$ は \emph{standard} と呼ばれる。
この二つの可能性と standard の定義は
\cite[命題5.3、定義5.4]{QuadricMU2021} による。
$X,Y$ は次数の移動を除いて $Q^{\oplus2}$ の部分加群となるので、
$g$ は両者の上で正則である
\cite[命題5.4、補題5.5]{SmithVdB}。

standard 条件が pencil の一般点で満たされることを、$E$ への制限で確認する。
一つの ruling の line module を $L_D=Q/V(D)Q$ とし、
$D$ は $E$ 上の次数2の有効因子、
$V(D)=H^0(E,\mathcal L(-D))$ とする。
ここでは右加群の規約を用いる。
この secant の記述は \cite[Basic Facts 7.2]{RSS2023} の右加群版による。
line module と spinor module の対応により
\[
 0\longrightarrow M(-1)\longrightarrow Q\longrightarrow L_D\longrightarrow0
\]
があり、$M$ は $X,Y$ のいずれかである。
\cite[命題7.2]{SmithVdB}
残る中心元 $g$ は $L_D$ 上で正則である。
実際 $L_D$ は2-critical であり、$gL_D=0$ なら
$B(E,\mathcal L,\alpha)$-module としての Hilbert 多項式の一次係数は
$4$ の非負整数倍になるが、line module の一次係数は1である。
従って $gL_D\ne0$ であり、criticality より $g$ の核は零である。
よって上の列は $E$ への制限後も完全である。
$\mathcal L(-D)$ は次数2で globally generated なので
\[
 Li^*M(-1)=\mathcal O_E(-D).
\]
$N=\mathcal O_E(D)$ とおく。
右加群の grading shift は
$F\mapsto\alpha_*(F\otimes\mathcal L)$ に対応するので
\[
 Li^*M=\alpha_*(\mathcal L\otimes N^{-1}).
\]
他方、\cite[補題5.5(2)]{SmithVdB} の列
\[
 0\longrightarrow M(-1)\longrightarrow Q^{\oplus2}
 \longrightarrow\Omega M(1)\longrightarrow0
\]
を制限する。$\Omega M(1)$ も $X,Y$ の一方であり、$g$ はその上で正則である。
従って制限後も完全で、行列式から $Li^*\Omega M(1)=N$ を得る。
従って $\Omega M(1)\cong M$ なら
\begin{equation}\label{blowup:nonstandard-necessary}
 N\otimes\alpha^*N\cong\mathcal L.
\end{equation}
写像 $\operatorname{Pic}^2E\to\operatorname{Pic}^4E$、
$N\mapsto N\otimes\alpha^*N$ は、平行移動を除けば multiplication by 2 である。
式\eqref{blowup:nonstandard-necessary}の解は高々4個である。
secant の因子類 $N\in\operatorname{Pic}^2E$ を固定すると、
対応する line module を annihilate する central quadric は一意である。
\cite[\S10.3]{SmithVdB} の左加群での記述に対し、
右加群の記述は \cite[Basic Facts 7.2]{RSS2023} に明記されている。
従って non-standard な $Q$ は、上の高々4個の $N$ が定める有限個の
pencil の元に含まれる。これらと非 smooth な元を除けば $Q$ は standard である。

この開集合上で、Mori--Ueyama の full geometric helix
\cite[定理5.16]{QuadricMU2021}と
$0\to X(-2)\to Q(-1)^{\oplus2}\to Y(-1)\to0$ を用いると
\[
 \bigl(\mathcal O_Q(-1),\widetilde Y(-1),\widetilde X(-1),\mathcal O_Q\bigr)
\]
という full strong exceptional collection が得られる。
強例外性と二つの中間対象の完全直交性は
同文献の補題5.6、5.9、5.13からも直接従う。
$E$ への制限は全て直線束で、その次数は順に $-4,-2,-2,0$ である。
従って最初の三頂点は再び hereditary fork になり、
補題\ref{blowup:quadric-kernel}の rank と degree の証明がそのまま適用できる。
残りの仮定は RSS の quadric elliptic surface の設定で満たされ、
$g$ は中心的な次数2の正則元である。
$Q$ は AS-Gorenstein で depth が3なので、
$Q_n=H^0(\operatorname{qgr}Q,\mathcal O_Q(n))$ for $n\ge0$ も成り立つ。
定理\ref{blowup:quadric-section-equivalence}の証明を適用して結論を得る。
\end{proof}

\begin{corollary}[共通切断環の二つの分解と三次元の積]
\label{blowup:rss-common-sod}
命題\ref{blowup:generic-sklyanin-quadric}の開集合から $Q$ を取り、
$\tau=\alpha^2$、$|\tau|=\infty$ とする。
任意の $t\in E$ に対して、RSS の同型
\cite[定理7.3]{RSS2023}を
\[
 R=Q^{(2)}(t)
   \cong S^{(3)}(\tau^{-1}p+\tau^{-1}q),
 \qquad \sigma^3=\tau,\qquad p\ne q
\]
と書く。ここで $S=S(E,\sigma)$ は三次元 Sklyanin 代数である。
$\mathcal A_U=\operatorname{Perf}(\operatorname{qgr}U)$ とおくと、
dg-enhanced な fully faithful functor
$\Phi_S\colon\mathcal A_S\to\mathcal A_R$、
$\Phi_Q\colon\mathcal A_Q\to\mathcal A_R$ と例外対象
$E_1,E_2,F\in\mathcal A_R$ が存在して
\begin{align}
 \mathcal A_R
  &=\langle E_2\otimes_k\operatorname{Perf}(k),\,
          E_1\otimes_k\operatorname{Perf}(k),\,
          \Phi_S(\mathcal A_S)\rangle,
       \label{blowup:rss-plane-sod}\\
 \mathcal A_R
  &=\langle F\otimes_k\operatorname{Perf}(k),\,
          \Phi_Q(\mathcal A_Q)\rangle.
       \label{blowup:rss-quadric-sod}
\end{align}
両 embedding は、RSS の環包含と上の同型が定める
\[
 Q_0(S)\cong Q_0(R)\cong Q_0(Q)
\]
という generic quotient の同一視に整合する。

さらに $\mathcal P=\operatorname{Perf}(\mathbb P^1)$ とし、
テンソル積は dg Morita テンソル積の冪等完備化とする。
$\mathcal B=\mathcal P\otimes_k^{\mathrm L}\mathcal A_R$ は
\begin{align}
 \mathcal B
   &=\langle\mathcal P,\,\mathcal P,\,
            \mathcal P\otimes_k^{\mathrm L}\mathcal A_S\rangle,
        \label{blowup:rss-product-plane-sod}\\
 \mathcal B
   &=\langle\mathcal P,\,
            \mathcal P\otimes_k^{\mathrm L}\mathcal A_Q\rangle.
        \label{blowup:rss-product-quadric-sod}
\end{align}
を持つ。各成分は上の fully faithful functor の積による像を表す。
三次元の generic quotient でも、$u$ を中心的な不定元として
$Q_0(S)(u)\cong Q_0(Q)(u)$ が保たれる。
\end{corollary}

\begin{proof}
quadric 側は命題\ref{blowup:generic-sklyanin-quadric}から従う。
平面側では $p_1=\tau^{-1}p$、$p_2=\tau^{-1}q$ とおく。
RSS の反復公式 \cite[命題5.4]{BlowupRSS2015} により
$S^{(3)}(p_1+p_2)=S^{(3)}(p_1)(p_2)$ である。
各段階の切断環は、上の次数1部分と $g$ による商の比較によって
Van den Bergh の切断環と一致する。
$p_1\ne p_2$ なので $\sigma p_1\ne\sigma p_2$ である。
二点の場合、これは \cite[定義11.1.2]{BlowupVdB2001} の一般位置条件を満たす。
従って同文献の定理11.1.3による二段階の比較同値と、
各段階の blowup 分解である定理8.4.1を合成すれば
式\eqref{blowup:rss-plane-sod}を得る。
$E_1$ は第一段階の例外対象を第二段階の fully faithful functor で移した対象であり、
$E_2$ は第二段階の例外対象である。
generic quotient での整合性は、各 Rees bimodule と切断環包含の
generic localization についての先の議論を二回適用すれば従う。

dg-enhanced な admissible decomposition の各 inclusion と adjoint を
$\mathcal P$ とテンソルすると、その単位・余単位および射影三角もテンソルされる。
従って完全忠実性、半直交性、生成性が保たれ、積の二つの分解を得る。
最後に $\mathbb P^1$ の generic point に制限し、
中心的な変数 $u$ を加えた Ore 分数環を取れば最後の主張を得る。
この系は固定された無限位数パラメータでの主張である。
\end{proof}
